% 本文件由 LaTeX 排版 Agent 根据有效 translation.json 自动生成。
% author=Miscellaneous; work_id=0514; title=驳斥诺瓦提安

\chapter{\WorkTitle{驳斥诺瓦提安}}

% structure: p0001 source=h1 target=omitted reason=page-title-already-rendered-by-parent

\Emph{不应拒绝失足者对宽恕的希望。}\par

1. 当我沉思并焦躁不安地在心中反复思索，该如何对待那些可悲的弟兄时——他们并非出于自愿，而是因狂怒的魔鬼突袭而受伤，至今仍已久经时日，在刑罚中忍耐——忽然，另一个仇敌出现在我面前，他与他自己的父爱相敌对，即异端者\Person{诺瓦提安}；他不仅如福音所预示的那样，像司铎或肋未人一样，从倒地的受伤者身旁经过，反而以巧妙而新颖的残酷方式，夺去救恩的希望、否认天父的怜悯、拒绝兄弟的悔改，意图杀死那受伤者。令人惊奇的是，有多少事物如此苦涩、严厉、乖谬！但人更容易看见别人眼中的草屑，却看不见自己眼中的横梁。然而，亲爱的弟兄们，不要被那背信弃义的异端者骤然狂热所动摇或困扰；虽然他陷于如此严重的分歧与分裂之罪中，并与教会分离，却以亵圣的狂妄，毫不畏惧地将他的指控掷回我们身上；因为虽然他自身因亵圣的污秽而变得不洁，他却争辩说我们才是那样。尽管经上记载狗应留在外面，宗徒也教导说这些狗必须躲避，正如我们所读，因为他说：\Quote{你们应提防狗，应提防作恶的工人} \BibleRef{斐 3:2}，他却不停止以吠叫激起他的疯狂，如同豺狼寻求幽暗之夜，以便在黑暗洞穴中轻易撕裂从牧人那里抢来的羊。当然，他声称他和那些被他聚集的朋友是金子。我们也毫不怀疑，那些离弃教会、成为背教者的人，如今容易变成金子；但那金子必定是以色列子民早期罪恶所象征的金子。但那些从埃及人手中夺来的金银器皿，仍处于主的大能之下，即基督的教会中；\Person{诺瓦提安}啊，如果你继续留在这房屋中，你可能也曾是一件宝贵的器皿；但现在你既不察觉，也不哀叹自己已变成了糠秕和稻草。\par

2. 那么，你为何因虚浮之事而自高自大？你将会遭受亏损而非益处。为何你因自己变得更贫穷而相信自己富有？请在默示录中聆听主的声音以正义的斥责责备你：\Quote{你说：}祂说，\Quote{我是富足的，已经发了财，什么都不缺；却不知道你是那困苦、可怜、贫穷、瞎眼、赤身的。} \BibleRef{默 3:17} 让那离弃基督的教会、以蒙蔽的理性毫不畏缩地转向那些分裂的鲁莽首领和纷争的作者（若望称之为假基督、福音作者比作糠秕、主基督称之为贼和强盗）的人，无论他是谁，都确信自己拥有这贫穷的财富。正如祂自己在福音中所宣告的，说：\Quote{那不从门进入羊栈，而从别处爬进去的，便是贼，是强盗。} \BibleRef{若 10:1} 此外，在同一处祂也说：\Quote{凡在我以先来的，都是贼和强盗。} \BibleRef{若 10:8} 这些人若不是离弃信仰、违犯天主教会、对抗天主命令的人，又是谁？圣神正当地藉先知斥责他们说：\Quote{你们筹划，却不由于我；\Emph{缔结了}盟约，却不由于我的圣神，以致罪上加罪。} \BibleRef{依 30:1} 那么，诺瓦提安那些最乖谬的朋友——即使现在也是极少数不幸的人——对此还能说什么？他们竟狂乱到如此愚蠢的地步，以致对天主或人都毫无敬畏。在他们中间，主教职被无耻地、毫无祝圣法规地追求着；但在我们中间，主教职在它自己的、由天主交付的宝座上被弃绝。那里真理说：\Quote{他们离弃我，却向我献祭；他们不献以色列子民圣洁的祭物，也不前来奉献至圣之物，却要在他们陷入的错误中蒙受羞辱。} \BibleRef{则 44:10} 只需寥寥数语就已足以证明他们是何等样人。因此，诺瓦提安人啊，请听：你们中间，属天的圣经被诵读，却未被理解；好吧，如果它们未被篡改的话。因为你们的耳朵被封闭，心灵被蒙蔽，以至于不接受任何属灵和救恩的光照；正如依撒意亚说：\Quote{天主的仆人们被蒙蔽了眼睛。} \BibleRef{依 13:19} 他们理当被蒙蔽，因为分裂者的愿望不在法律之内；这法律指示我们那唯一教会就在那只方舟里——这方舟是洪水之前，在天主眷顾下由诺厄建造的；方舟中不仅容纳洁净的动物，也容纳不洁净的——诺瓦提安啊，快速回答你，我们发现那只方舟独自分担了得救，连同其中所有的一切，而其他不在方舟内的事物都在洪水中灭亡了。从那只方舟中放出了两只鸟，一只乌鸦和一只鸽子；这乌鸦确实代表了不洁之人的形象或预表，这些人将借着世界的宽广大道永远处于黑暗之中，也代表了将要兴起的背教者，他们以不洁之物为食，最终不转向教会；正如我们所读，我们发现乌鸦被放出后，再也没有回来。凡是与这只鸟相似的人，即不洁之灵，将再也不能回到教会，因为主必会禁止他们，即使他们愿意，正如祂命令梅瑟说：\Quote{凡患麻风病和不洁的，都要赶出营外。} \BibleRef{户 5:2} 但那只放出去后返回的鸽子，象征着那毫不迟延的人，因为他的脚找不到安息。诺厄把它接入方舟；第七天再把它放出去时，又接了回来，它嘴里衔着一片橄榄叶。\par

3. 亲爱的弟兄们，我并非轻率地思考这些事，也不是出于人的智慧，而是因天上之主的俯就，允许我们的心灵必要而恰当地领受——我说，那只鸽子本身向我们表明一种双重类型。从前，即从神圣管理之始，它便暗示了其自身的形象，首先是首要的——即圣神的形象。藉着它的口，为人类救恩而准备的圣洗圣事，并依照天上的计划，仅在教会内举行。此外，它三次从方舟被放出，在水面上空飞翔，已经预示了我们教会的圣事。因此，主基督也嘱托伯多禄，以及其余的门徒说：\Quote{你们要去向万民宣讲福音，因父、及子、及圣神之名给他们授洗。} 也就是说，那在诺厄时代藉着鸽子以预像方式运作的同一圣三，如今在教会内藉着门徒以属神的方式运作。\par

4. 现在我们也来看那从方舟放出的鸽子的第二个特征，就是说，在洪水时期，当所有深渊迸裂，天上的水闸向大地敞开，因人在主面前日日行恶；正如梅瑟所说：\Quote{上主天主见人在地上的罪恶很大，他们从幼年以来，终日心中思念的尽是邪恶；祂说：我要将我所造的人从地面上消灭，从人到牲畜，到爬虫，到天空的飞鸟。} \BibleRef{创 6:5} 因此，在洪水时期，鸽子从方舟被放出，那时大水以全力猛烈冲击大地。\par

5. 那方舟承载着教会的预像，正如我们上面所说，它被汹涌的洪水猛烈冲击。因此，诺厄时代的洪水预示了最近波及全世界的迫害。再者，那从四面八方汇合、不断上涨的洪流，象征了那些为毁灭教会而聚集的人民，正如\WorkTitle{默示录}教导说：\Quote{你所看见的水，是指诸民族、诸民众、诸邦国。} \BibleRef{默 17:15} 此外，那只找不到落脚之处的鸽子，代表了跌倒者的形象，他们因忘却天主的宣告而跌倒，或是因无知而单纯，或是因狂妄而伪装。主曾在福音中预言他们未来的毁灭，说：\Quote{凡听了我这些话而不实行的，就好像一个愚昧人，把房子盖在沙土上：雨淋、水冲、风吹，袭击那座房子，它就倒塌了，且倒塌得很厉害。} \BibleRef{玛 7:26} 为了不让我们显得轻率地以那只鸽子作为跌倒者的形象，先知责备那座城如同一只鸽子，即跌倒者的特征，说：\Quote{鸽子不听声音；即那闻名而被赎的城不接受教训，也不信赖上主。}\par

6. 此外，那只鸽子找不到落脚之处，如上所述，这象征了那些背弃者的脚步；即那些被发光的毒蛇咬伤、祭献而跌倒的人，他们不能再踩在毒蛇和蜥蜴上，也不能践踏龙和狮子。因为主曾将这权柄赐给祂的门徒，正如祂在福音中所说：\Quote{看，我已给你们权柄，可践踏一切仇敌的势力，并能踏在蛇和蝎子上，它们决不会伤害你们。} \BibleRef{路 10:19} 因此，当这么多恶灵为毁灭跌倒者而攻击、鼓动时，为受伤者提供了一条得救之路，使他们能竭尽全力拖着全身，回到自己的营寨，在那里被收容后，以属神的药物医治他们的伤口。这样，鸽子在隔了几天后被再次放出，返回时不仅显示了坚定的脚步，而且口中衔的橄榄枝还显示了和平与胜利的标志。因此，这两次放出向我们表明了双重的迫害考验：第一次，他们中跌倒的人成了失败者；第二次，他们中跌倒的人却成了得胜者。因为亲爱的弟兄们，我们中无人怀疑或不确定，那些在第一次斗争——即德西乌斯迫害——中受伤的人，后来在第二次交锋中如此勇敢地坚持，以至于蔑视世俗统治者的谕令，保持了不可征服的姿态；因为他们不怕，效法善牧的榜样，舍弃生命，流血牺牲，不躲避暴君的任何野蛮行径。\par

7. 看，这些分裂分子毫不畏惧地称之为\Quote{木、草、禾秸} \BibleRef{格前 3:12} 的人民，在主眼中是何等光荣、何等可亲！与他们同等的人，即那些仍处于同样跌倒罪过中的人，他们竟以为绝不可允许他们悔改。\Emph{他们如此判断}，是出于主的那句话，祂说：\Quote{凡在人面前否认我的，我在天父面前也必否认他。} \BibleRef{玛 10:33} 呜呼哀哉！他们为何反抗主的诫命，使这诺瓦提安的子孙，效法其父魔鬼，竟试图执行基督要在祂审判之时才做的事？\Emph{即}，当\WorkTitle{圣经}说：\Quote{复仇属于我，我必报应，主说。} \BibleRef{希 10:30}\par

8. 我们将回答他们\Emph{关于}主所说的那番话，他们误解了它，也向自己解释错了。因为祂说：\BibleQuote{凡在人前否认我的，我在我天之圣父面前也必否认他}，这意思确实是关乎将来的时期——就是主要开始审判人隐密事的时期——就是我们必须都站在基督审判台前的时期——就是许多人要开始说：\BibleQuote{主啊，主啊，我们不是因祢的名预言，因祢的名驱魔，因祢的名行了许多奇迹吗？}\BibleRef{玛 7:22}然而\Emph{他们}却要听见主的声音说：\BibleQuote{离开我，你们这些作恶的人：我不认识你们。}\BibleRef{玛 7:22}那时，祂所说的\BibleQuote{我也必否认他们}就要应验。但主基督主要否认谁呢？岂不是否认你们所有的异端者、分裂者以及与祂的名无关的人吗？因为你们这些一度是基督徒、如今却成了诺瓦蒂安派的人，由于你们名称的称呼，用后来的背信改变了你们最初的信仰。我希望你们回答自己的主张。请阅读并教导：我们主曾否认哪些失败或否认祂的人，当祂仍与他们在一起时？然而，对于其他仍与祂在一起的门徒，祂说：\BibleQuote{你们也要离开吗？}\BibleRef{若 6:67}甚至\Person{伯多禄}，祂曾预言他将要否认祂，当他否认了祂时，祂并没有否认他，而是扶持了他；并且当后来他痛苦地哀哭自己的否认时，主亲自安慰了他。\par

9. 你们是多么愚蠢，\Person{诺瓦蒂安}，只读那些导致救恩毁灭的经文，却忽略那些导向怜悯的经文，而圣经呼喊说：\BibleQuote{悔改吧，你们这些迷途的人：在心上归正。}\BibleRef{则 18:30}同一先知也劝勉说：\BibleQuote{你们要全心归向我，禁食、哭泣、哀悼；要撕裂你们的心，而不是衣服；归向你们的天主，因为祂是仁慈的，且以大怜悯怜恤人？}\BibleRef{岳 2:12}\par

10. 因此我们听见主有大怜悯。让我们听一听圣神藉达味所作的见证：\BibleQuote{如果他的子女离弃我的法律，不遵行我的诫命；如果他们亵渎我的公义，不守住我的规条，我必用杖责罚他们的过犯，用鞭责罚他们的罪恶。但我必不将我的慈爱尽行收回。}类似的话我们也读到主藉厄则克耳所说：\BibleQuote{人子，以色列家族居住在自己的土地上，他们以自己的罪恶玷污了那地：他们的不洁在我面前如同行经妇人的不洁。我向他们倾泻我的愤怒，将他们分散在列国中；我按他们的罪恶审判了他们，因为他们玷污了我的圣名；并且论到他们，曾说：\InnerQuote{这是主的百姓，我因我的圣名饶恕了他们，因为以色列家族在列邦中藐视了我的圣名。}}\BibleRef{则 36:17}与此相关，祂说：\BibleQuote{因此你要对以色列家族说：主这样说：我绝不饶恕你们，以色列家族；但我要因我的圣名饶恕你们，你们在列邦中玷污了我的圣名；当我在你们中被圣化时，你们就知道我是主。}主也向同一人说了同样的话：\BibleQuote{人子，你要对以色列家族说：你们为什么说话，说：我们在罪恶中衰败，我们怎么能得救？你要对他们说：我指着我的永生说，主说：我断不喜悦罪人死去；却喜悦罪人离开他的恶道，得以存活。所以你们要离开你们的恶道：以色列家族，你们为什么将自己交给死亡？}\BibleRef{则 33:10}同样，依撒意亚先知说：\BibleQuote{我不会永远向你们发怒，也不会始终不保护你们。}\BibleRef{依 57:16}并且因为耶肋米亚先知以罪恶百姓的身份向主祈祷说：\BibleQuote{主啊，求祢在审判中改正我们，不要在发怒时，免得祢使我们人数稀少。}\BibleRef{耶 10:24}依撒意亚也接着说：\BibleQuote{因他的罪，我稍微磨难了他，击打了他，并转脸不看他；他受了磨难，忧苦地走他的路。}\BibleRef{依 57:17}并且因为他劳苦，他接着说：\BibleQuote{我看到了他的道路，医好了他，给了他真正的劝勉，平安又平安。}\BibleRef{依 57:19}为的是那些悔改、祈祷并劳苦的人，复原是可能的，因为他们会悲惨地丧亡，并会背离基督。\par

11. 此外，这也在福音中得到证明，那里描述了一个罪妇，她来到一个法利塞人的家里，主和祂的门徒应约赴宴；她带着一玉瓶香膏，站在主身后，用眼泪洗祂的脚，用头发擦干，并连连亲祂的脚；以致那法利塞人被激怒，说：\BibleQuote{如果这人是先知，必定知道摸祂的是谁，是怎样的女人；因为她是一个罪妇。}于是主立即，这位赦罪者和忏悔者的接纳者，说：\BibleQuote{西满，我有话要对你说。他回答说：夫子，请说。主说：有一个债主有两个债务人：一个欠五百银钱，另一个欠五十。因为他们无力偿还，债主都免了他们的债。祂问：这两人中哪一个爱得更多？西满回答说：当然是被赦免更多的那个。祂接着说：你看见这妇人吗？我进了你的家，你没有给我亲嘴；但这妇人从我来时起，就不住地亲我的脚。你没有用水洗我的脚，但这妇人用眼泪洗了我的脚，用头发擦干。你没有用油抹我的脚，但这妇人用香膏抹了我的脚。所以我告诉你，西满，她的罪都赦免了。}看，主以其慷慨的仁慈免了两人的债；看，祂赦免罪过；看，那罪妇忏悔、哭泣、祈祷，并获得了罪的赦免！\par

12. 如今，你若能，就脸红吧，\Person{诺瓦齐安}；停止以你不敬的论调欺骗无知者；停止以某一点的微妙之处惊吓他们。我们诵读、钦崇，并不越过主那属天的审判，祂说：凡否认祂的，祂也要否认他。但这指的是忏悔者吗？我何苦要费时举证个别的仁慈事例呢？因为天主的仁慈并未拒绝尼尼微人，他们虽是外邦人，远离主的律法，但因城邑即将倾覆而恳求仁慈。也没有拒绝法郎本人，他虽以渎神的狂妄抗拒，先前遭受从天而降的灾祸，便转向梅瑟和他的兄弟说：\Quote{请为我祈求上主，因为我犯了罪。} \BibleRef{出 9:28}天主的愤怒立刻从他身上止息了。而你，噢，\Person{诺瓦齐安}，竟论断并宣告失足者无望获得平安和仁慈，也不肯垂听宗徒的斥责，他说：你是谁，竟论断别人的仆人？他或站或跌，自有他的主人。而且他必站得住，因为天主有能力使他站住。\BibleRef{罗 14:4}因此，圣神合宜而必要地藉着那些失足者之口斥责你，祂说：\Quote{我的仇敌啊，不要向我夸耀；因为我虽跌倒，却要起来；我虽坐在黑暗里，上主是我的光明。我要忍受上主的盛怒，因为我得罪了祂，直到祂为我辩护，并为我施行审判和正义，领我进入光明。我必得见祂的公义；我的仇敌看见，就必抱愧蒙羞。} \BibleRef{米 7:8}\par

13. 我恳求你，难道你没有读过：\Quote{不要自夸，不要口出狂言，不要让傲慢从你口中发出；因为上主从地上提拔穷乏人，从粪堆高举乞丐，使他们与民间的首领同坐？} 难道你没有读过：\Quote{上主抵挡骄傲的人，赐恩宠给谦卑的人？} \BibleRef{雅 4:6}难道你没有读过：\Quote{凡自高的，必降为卑？} \BibleRef{玛 23:12}难道你没有读过：\Quote{天主消灭骄傲人的纪念，却不忘记谦卑人的回忆？} 难道你没有读过：\Quote{你们用什么判断来判断，你们也要受什么判断？} \BibleRef{玛 7:2}难道你没有读过：\Quote{凡恨自己弟兄的，就是在黑暗里，且在黑暗里行，也不知道往哪里去，因为黑暗弄瞎了他的眼睛？} \BibleRef{若一 2:11}那么，这\Person{诺瓦齐安}为何变得如此邪恶、如此迷失、如此疯狂地充满纷争的怒火，我无法明白，因为他原本常在同一家中——即基督的教会——为自己邻人的罪孽哀哭，如同为自己的罪一样；他本会背负弟兄的重担，正如宗徒所劝勉的；他本会用属天的劝言坚固那些在信德上动摇的人。但现在，自从他开始推行那仅以杀人为乐的加音异端以来，他甚至近来也不肯饶恕自己。但他若读过\Quote{义人的义在他犯罪之日不能救他，恶人的恶在他悔改之日也不能害他，} \BibleRef{则 33:12}，他早就披麻蒙灰忏悔了；他却一贯反对忏悔者；他更乐意毁坏已建立和屹立的事物，而不愿重建倾颓的事物；他因错误的反对恐吓了我们许多最可怜的弟兄，使他们再次成为外教人，声称失足者的忏悔是徒劳的，不能为他们获得救恩，尽管圣经大声疾呼说：\Quote{你要回想你是从哪里跌落的，并要悔改，否则我要临到你，除非你悔改。} \BibleRef{默 2:5}确实，当圣经写信给七个教会，责备她们各人的罪行和罪恶时，祂说：悔改。除了祂用自己宝血重价赎回的人之外，这些还能是对谁说的呢，无疑地？\par

14. 噢，你这不敬和邪恶的异端者\Person{诺瓦齐安}！你在过去的日子里曾知道教会内有许多重大罪行是自愿犯下的，而你自己在成为天主家中的背教者之前，确实教导过：犯罪者若随后行善，这些罪可从记忆中被抹去；按照圣经的信德，经上说：\Quote{若恶人转离他所犯的一切罪恶，且遵行法律，行正直合理的事，他必得永生，且不因他的恶行而死。} \BibleRef{则 18:21}因为他所犯的罪，必因随后的善行而从记忆中被抹去。现在请再思，那些失足跌倒、被魔鬼剥得赤露的人，他们的创伤是否应当医治？他们既已\Emph{如被击倒}，被\Quote{那蛇从口中向妇人吐出的洪水般的暴力} \BibleRef{默 12:15}所打击。但宗徒说：\Quote{我该说什么呢？} \Quote{我称赞你们吗？在这件事上我不称赞，因为你们聚在一起，不是为得益，而是为受损。} \BibleRef{格前 11:17}因为哪里有\Quote{嫉妒纷争，你们岂不是属肉体，照着人的样子而行吗？} \BibleRef{格前 3:3}我们实在不必奇怪，这\Person{诺瓦齐安}如今竟敢对失足者行如此邪恶、如此严厉的事，因为我们先前已有此类背信的例证。\Person{撒乌耳}，那个\Emph{曾经良善}的人（见\BibleRef{撒上 9:2}），除其他事外，后来被嫉妒所胜，竭力对\Person{达味}行一切严苛敌对之事。那个\Person{犹达斯}，被选在宗徒之中，在天主家中始终同心并忠心，后来他自己却背叛了天主。\par

主的确曾预言，将有许多披着羊皮的凶暴豺狼来到。这些凶暴豺狼，岂不就是那些以诡谲的意图阴谋摧残基督羊群的人吗？正如我们在\BibleBook{匝加利亚}{书}中所读到写的：\Quote{看，我要在土地上兴起一个牧人，他不眷顾离散的，却要吃肥羊的肉，撕裂它们的蹄。}\BibleRef{匝 11:16}同样，在\BibleBook{厄则克耳}{书}中，祂也斥责这类牧人，即盗贼和屠夫（我将照祂所想的来说），说：\Quote{牧人们，你们为何喝牛奶，吃奶饼，却使强壮的归于无有？你们不探望软弱的，不医治跛足的，不领回迷失的，竟让我的百姓在荆棘和蒺藜中流荡？为此，上主说：看，我要反对那些牧人，从他们手中追讨我的羊；我要驱逐他们，使他们不再牧放我的羊；我的羊不再成为他们的掠物。我要像牧人在有云和黑暗的日子寻找他的羊群一样，寻找我的羊。我要寻找我的羊，在他们所散居的各个地方寻找他们。我要寻找那失落的，领回那迷失的，医治那跛足的，看顾那软弱的；我要以公义牧放我的羊。}\EditorialNote{厄则克耳第三十四章}\par

15. 说这些话的是谁呢？当然是把那撇下九十九只羊，去寻找那一只从祂羊群中迷失的羊的那一位；正如达味所说：\Quote{我迷了路，像一只迷失的羊，}基督将它寻回，背在肩上，带回那柔弱的罪人；祂欢欣踊跃，召来朋友和家人，说：\Quote{与我一同欢乐罢，因为我那迷失的羊找到了。我告诉你们，}祂说，\Quote{对于一个悔改的罪人，天上也要这样欢乐。}\BibleRef{路 15:6}祂接着说：\Quote{或者一个妇女，有十个\OriginalTerm{银币}{denarius}，若失落一个，岂不点上灯，打扫房屋，细心寻找，直到找着吗？找着了，就请朋友邻人来，说：与我一同欢乐罢，因为我失落的那个\OriginalTerm{银币}{denarius}已经找到了。我告诉你们，对于一个悔改的罪人，在天主的使者面前也是这样欢乐。}\BibleRef{路 15:6}但另一方面，那些不为自己的恶行悔改的人，当从主自己的回答中知道等待他们的是什么；因为我们在福音中读到，\Quote{有些人从加里肋亚来到主那里，告诉祂关于比拉多使他们的血搀和在祭品中的事；主回答说：你们以为那些加里肋亚人比别的加里肋亚人更有罪，才遭受这事吗？不是的；我告诉你们，除非你们悔改，你们都要同样丧亡。或者那十八个人，史罗亚塔倒塌压死了他们，你们以为他们比所有住在耶路撒冷的人更负罪债吗？不是的；我告诉你们，}祂说，\Quote{除非你们悔改，你们都要同样丧亡。}\BibleRef{路 13:1}\par

16. 那么，亲爱的弟兄们，让我们尽量奋发起来；打破懒惰和安逸的沉睡，让我们警醒遵守主的诫命。让我们全心寻找我们所失去的，好能找着；因为经上说：\Quote{寻找的，就必得到；敲门的，就必给他开门。}\BibleRef{路 11:10}让我们以神灵的干净打扫我们的家，使我们心中每一秘密和隐藏的地方，真正被福音的光照亮，可以说：\Quote{唯独得罪了祢，在祢眼前行了这恶事。}因为罪人的死亡是邪恶的，在地狱里没有悔改。我们要特别默想审判和报应的日子，以及我们众人所必须相信并坚定持守的：\Quote{天主不看人的情面。}\BibleRef{羅 2:11}因为祂在\BibleBook{}{申命纪}中命令，审判时不可看人的情面：\Quote{不可看人的情面，}祂说，\Quote{也不可按最小或最大的来审判。}\BibleRef{申 1:17}祂也藉\BibleBook{}{厄则克耳}说了类似的话：\Quote{所有的灵魂，}祂说，\Quote{都属于我；父亲的灵魂，儿子的灵魂也一样：犯罪的灵魂，必要死亡。}\BibleRef{則 18:4}因此，我们必须敬畏祂，紧抓住祂，以我们圆满而相称的告解来恳求祂平息愤怒，\Quote{祂有权把灵魂和身体投入\OriginalTerm{革赫纳}{Gehenna}的火里，}\BibleRef{瑪 10:28}——正如所记载的：\Quote{看，祂率领祂的千万天使来临，要在众人身上行审判，消灭一切恶人，定一切血肉之躯的罪，因为恶人一切邪恶行为，以及罪人对天主所说的一切亵渎言语。}\par

17. 与此类似，达尼尔也说：\Quote{我观望，直到安置了宝座，万古常存者坐上去，祂的衣服洁白如雪，祂的头发洁白如羊毛，祂的宝座是火焰，祂的车轮是烈火。从宝座前有火河涌出，千万侍从服事祂，千万的侍者侍立在祂面前，审判者已坐堂，案卷已展开。}\BibleRef{达 7:9} 若望更清楚地论及审判之日和世界的终结，说：\Quote{当祂开启第六个印时，我看见，发生了大地震，太阳变黑像粗毛衣，整个月亮变成血；天上的星辰坠落到地上，如同无花果树被大风摇落未熟的果子。天也像是书卷被卷起来，一切山岭和岛屿都移动了位置。世上的君王、首领、将军、富人、壮士，以及所有为奴的、自由的，都隐藏在山洞和山穴中，向山岭和岩石说：\InnerQuote{倒在我们身上，隐藏我们，好躲避坐在宝座上的圣父的面容和羔羊的义怒，因为祂发怒的大日子来到了，谁能站立得住？}}\BibleRef{默 6:12} 若望在同一部《默示录》中也说，这也是启示给他的：\Quote{我看见了一个大白宝座，坐在上面的圣父，天地都从祂面前逃避，再也找不到它们的地方了。我又看见死过的人，无论大小，都站在宝座前，案卷被展开了；还有另一本书，即生命册，也被展开了；死过的人都按着他们行为所记载的，受审判。}\BibleRef{默 20:11} 此外，宗徒也给予善劝，如此劝诫我们说：\Quote{不要让任何人用空言欺骗你们；就是因为这些事，天主的义怒才降在悖逆之子身上。所以你们不要与他们同流合污。}\BibleRef{弗 5:6}\par

18. 因此，让我们以信德的全部力量赞美天主；让我们做完全的告解，因为诸天之上因我们的悔改而欢乐，众天使喜乐，基督也喜乐；祂再次以圆满而慈悲的温和劝诫我们，我们身负众罪、恶贯满盈，应当停止作恶，说：\Quote{转身，离开你们的恶行，你们的罪孽就不会成为你们的惩罚。把你们所犯的诸般恶行从你们身上抛弃，为自己造一颗新的心和一个新的精神。以色列家！你们为什么要自招死亡？因为我不喜欢丧亡者的死亡。}\BibleRef{则 18:30} \Quote{是我，是我，是我涂抹了你的过犯，我不再记念你的罪恶。你要提醒我，让我们一起辩论；你先提出你的案件来，好证明你无罪。}\BibleRef{依 43:25} 弟兄们，慈悲之路已敞开，让我们用圆满的补赎恳求天主；让我们自卑，好被高举；让我们顺服于神圣的劝诫，好能逃脱主的日子和祂的义怒。因为祂这样说：\Quote{我儿，你观看万民，要认识：谁曾信赖上主，而受辱？谁曾持守祂的诫命，而被遗弃？谁曾呼号祂，而祂未曾俯听？因为上主是慈悲仁爱的，在患难之日，赦免所有诚心寻求祂的人的罪过。}\BibleRef{德 2:10} 因此祂说：\Quote{先提出你的案件来，好证明你无罪。} 愿那充满告解（\OriginalTerm{告解}{confessio}）的祈祷先在你手中。\par

\SourceRecord{Translated by Robert Ernest Wallis.；Ante-Nicene Fathers；Vol. 5.；Edited by Alexander Roberts, James Donaldson, and A. Cleveland Coxe.；Buffalo, NY: Christian Literature Publishing Co.,；1886.；<http://www.newadvent.org/fathers/0514.htm>.}
